{\color{nuevo}% >>> CONTENIDO NUEVO (entrega 50%) <<<
\chapter{Metodología de desarrollo}
\label{cap:diseno}

Este capítulo describe cómo se construye la solución a partir del análisis de requerimientos del Capítulo~\ref{cap:requerimientos}. Se presenta la metodología de trabajo adoptada, la arquitectura de la solución y su infraestructura de red, las tecnologías seleccionadas y su justificación, el modelo de datos y, finalmente, la estrategia de validación del sistema.

% ═════════════════════════════════════════════════════════════════════════════
\section{Metodología}

El desarrollo del prototipo sigue un enfoque \textbf{iterativo e incremental}: la funcionalidad se construye en hitos sucesivos, cada uno con valor demostrable, priorizando primero la base transversal (autenticación, persistencia y ciclo de vida del envío) y avanzando luego hacia los componentes diferenciales (asignación inteligente, clasificación por visión computacional y cálculo de impacto ambiental). Este orden permite validar de forma temprana las decisiones de diseño fundamentadas en los antecedentes y ajustar el alcance según los resultados.

El trabajo se apoya en un conjunto de prácticas de ingeniería: control de versiones con Git y GitHub, empaquetado reproducible mediante contenedores, configuración externalizada por entorno, pruebas automatizadas sobre los módulos críticos e integración continua. La construcción del conjunto de datos para el clasificador se inicia en paralelo desde las primeras etapas, por ser la tarea que no puede acelerarse con código.

% ═════════════════════════════════════════════════════════════════════════════
\section{Arquitectura de la solución}

La solución adopta el patrón de \textbf{monolito modular orientado a dominio}: una única aplicación de servidor organizada en módulos autocontenidos, cada uno responsable de un dominio del negocio. Se descarta una arquitectura de microservicios porque su complejidad operativa (comunicación de red entre servicios, observabilidad distribuida y consistencia eventual) no se justifica en un prototipo de estas dimensiones; la modularización interna preserva, no obstante, límites claros que habilitarían la extracción de un servicio en el futuro si la escala lo requiriera.

La arquitectura se describe siguiendo el modelo \textbf{C4}, de mayor a menor nivel de abstracción.

\begin{itemize}[leftmargin=2em]
    \item \textbf{Contexto.} El sistema atiende a cinco tipos de actor (cliente particular, transportista, PyME fletera, PyME comercio y administrador) a través de dos frentes —una aplicación móvil para clientes y transportistas y un panel web para PyMEs y administración— sobre un único backend. Sus dependencias externas son mínimas y reemplazables: un servicio de ruteo (con alternativa de cálculo propia), un servicio de mapas y geocodificación, y la persistencia de datos e imágenes.
    \item \textbf{Contenedores.} Las unidades desplegables por separado son: la aplicación móvil, el panel web, la API REST, la base de datos y el almacenamiento de objetos. Las aplicaciones cliente se comunican con la API sobre HTTPS con autenticación por token; la API concentra toda la lógica de negocio y el acceso a datos.
    \item \textbf{Componentes.} Dentro de la API, cada módulo de dominio respeta una separación estricta en cuatro capas —enrutamiento (HTTP y validación), servicio (lógica de negocio), repositorio (acceso a datos) y esquemas—, con una regla de dependencia unidireccional. Los módulos implementados abarcan autenticación, usuarios, transportistas, envíos, trayectos, asignación, visión, seguimiento, capacidad, impacto ambiental, organizaciones y administración.
\end{itemize}

La Figura~\ref{fig:c4-contexto} representa el nivel de contexto: el sistema como una caja negra, sus usuarios y sus dependencias externas.

\begin{figure}[ht]
    \centering\color{nuevo}
    \begin{tikzpicture}
        \node[c4system] (S) at (0,0) {\textbf{DePaso}\\ Logística urbana\\ colaborativa};
        % Actores
        \node[c4person] (cli) at (-5.6,2.7)  {Cliente\\ particular};
        \node[c4person] (tra) at (-5.6,0.9)  {Transportista};
        \node[c4person] (pf)  at (-5.6,-0.9) {PyME fletera};
        \node[c4person] (pc)  at (-5.6,-2.7) {PyME comercio};
        \node[c4person] (adm) at (0,-3.4)    {Administrador};
        % Dependencias externas
        \node[c4ext] (osrm) at (5.6,1.8)  {Servicio de\\ ruteo};
        \node[c4ext] (maps) at (5.6,0)    {Mapas /\\ geocodificación};
        \node[c4db]  (db)   at (5.6,-2.1) {Base de datos\\ + almacenamiento};
        % Actores -> sistema
        \draw[c4arrow] (cli) -- (S);
        \draw[c4arrow] (tra) -- (S);
        \draw[c4arrow] (pf)  -- (S);
        \draw[c4arrow] (pc)  -- (S);
        \draw[c4arrow] (adm) -- (S);
        % Sistema -> externos
        \draw[c4arrow] (S) -- (osrm) node[c4lbl] {distancias};
        \draw[c4arrow] (S) -- (maps) node[c4lbl] {direcciones};
        \draw[c4arrow] (S) -- (db)   node[c4lbl] {persistencia};
    \end{tikzpicture}
    \caption{Diagrama de contexto del sistema (C4, nivel 1)}
    \label{fig:c4-contexto}
\end{figure}

La Figura~\ref{fig:c4-contenedores} desglosa el sistema en sus contenedores desplegables.

\begin{figure}[ht]
    \centering\color{nuevo}
    \begin{tikzpicture}
        \node[c4container] (app) at (-2.9,3.2) {App móvil\\ (clientes y\\ transportistas)};
        \node[c4container] (web) at (2.9,3.2)  {Panel web\\ (PyMEs y\\ administración)};
        \node[c4container, minimum width=3.4cm] (api) at (0,0.6) {\textbf{API REST}\\ lógica de negocio};
        \node[c4ext] (osrm) at (5,0.6)     {Servicio de\\ ruteo};
        \node[c4db]  (db)   at (-2.9,-2.6) {Base de\\ datos};
        \node[c4db]  (st)   at (0.2,-2.6)  {Almacenamiento\\ de objetos};
        \node[c4ext] (ml)   at (3.2,-2.6)  {Modelo de\\ IA (.\,keras)};
        \draw[c4arrow] (app) -- (api) node[c4lbl] {HTTPS + \emph{token}};
        \draw[c4arrow] (web) -- (api) node[c4lbl] {HTTPS + \emph{token}};
        \draw[c4arrow] (api) -- (db);
        \draw[c4arrow] (api) -- (st);
        \draw[c4arrow, dashed] (api) -- (ml)   node[c4lbl] {carga inicial};
        \draw[c4arrow, dashed] (api) -- (osrm) node[c4lbl] {o alternativa};
    \end{tikzpicture}
    \caption{Diagrama de contenedores (C4, nivel 2)}
    \label{fig:c4-contenedores}
\end{figure}

Internamente, cada contenedor de la API se organiza en módulos de dominio con cuatro capas, como muestra la Figura~\ref{fig:c4-componentes}.

\begin{figure}[ht]
    \centering\color{nuevo}
    \begin{tikzpicture}
        \node[c4container] (l1) at (0,2.4)  {Enrutamiento\\ (HTTP, validación)};
        \node[c4container] (l2) at (0,0.8)  {Servicios\\ (lógica de negocio)};
        \node[c4container] (l3) at (0,-0.8) {Repositorios\\ (acceso a datos)};
        \draw[c4arrow] (l1) -- (l2);
        \draw[c4arrow] (l2) -- (l3);
        \node[c4ext, text width=3cm] (mods) at (-4.3,0.8) {Módulos de dominio: auth, usuarios, transportistas, envíos, \emph{matching}, trayectos, seguimiento, visión, CO\textsubscript{2}, organizaciones, administración};
        \node[c4ext, text width=2.6cm] (core) at (4.3,2.2) {\emph{core}: configuración, seguridad, base de datos};
        \node[c4ext, text width=2.6cm] (shr)  at (4.3,-0.6) {\emph{shared}: geo, ruteo, enums, repositorio base};
        \node[c4db] (db) at (0,-2.7) {Base de datos};
        \draw[c4arrow] (l3) -- (db);
        \draw[c4arrow, dashed] (mods) -- (l1);
        \draw[c4arrow, dashed] (l1) -- (core);
        \draw[c4arrow, dashed] (l2) -- (shr);
        \node[boundary, fit=(l1)(l2)(l3), inner sep=5mm] {};
    \end{tikzpicture}
    \caption{Componentes del backend: módulos de dominio en cuatro capas (C4, nivel 3)}
    \label{fig:c4-componentes}
\end{figure}

Finalmente, la Figura~\ref{fig:secuencia} integra los tres pilares del proyecto —clasificación, asignación y cálculo de impacto— en el flujo de un envío colaborativo de punta a punta.

\begin{figure}[ht]
    \centering\color{nuevo}
    \resizebox{\textwidth}{!}{%
    \begin{tikzpicture}[
        msg/.style={-{Stealth[length=2mm]}, font=\scriptsize},
        ret/.style={-{Stealth[length=2mm]}, dashed, font=\scriptsize},
        part/.style={draw, fill=RoyalBlue!12, minimum height=0.8cm, font=\footnotesize, align=center},
    ]
        \node[part] (C) at (0,0)    {Cliente};
        \node[part] (A) at (3.2,0)  {App móvil};
        \node[part] (P) at (6.8,0)  {API};
        \node[part] (D) at (10.2,0) {Base de datos};
        \node[part] (R) at (13.2,0) {Cadete};
        \foreach \n in {C,A,P,D,R} { \draw[dashed, black!45] (\n.south) -- ($(\n.south)+(0,-12.8)$); }
        \draw[msg] (C|-0,-1.2)  -- (A|-0,-1.2)  node[midway,above]{fotografía el paquete};
        \draw[msg] (A|-0,-2.2)  -- (P|-0,-2.2)  node[midway,above]{clasificar carga};
        \draw[ret] (P|-0,-3.2)  -- (A|-0,-3.2)  node[midway,above]{categoría + confianza};
        \draw[msg] (A|-0,-4.2)  -- (P|-0,-4.2)  node[midway,above]{cotizar (dedicado / colaborativo + CO\textsubscript{2})};
        \draw[msg] (A|-0,-5.2)  -- (P|-0,-5.2)  node[midway,above]{crear y pagar el envío};
        \draw[msg] (P|-0,-6.2)  -- (D|-0,-6.2)  node[midway,above]{registrar envío y pago (garantía)};
        \draw[msg] (R|-0,-7.2)  -- (P|-0,-7.2)  node[midway,above]{consultar pedidos compatibles};
        \draw[msg] (R|-0,-8.2)  -- (P|-0,-8.2)  node[midway,above]{aceptar el envío};
        \draw[ret] (P|-0,-9.2)  -- (C|-0,-9.2)  node[midway,above]{notificar asignación};
        \draw[msg] (R|-0,-10.2) -- (P|-0,-10.2) node[midway,above]{actualizar estado (retiro $\to$ tránsito $\to$ entregado)};
        \draw[msg] (P|-0,-11.2) -- (D|-0,-11.2) node[midway,above]{liberar pago y persistir CO\textsubscript{2}};
        \draw[msg] (C|-0,-12.4) -- (P|-0,-12.4) node[midway,above]{calificar};
    \end{tikzpicture}%
    }
    \caption{Diagrama de secuencia: envío colaborativo de punta a punta}
    \label{fig:secuencia}
\end{figure}

% ─────────────────────────────────────────────────────────────────────────────
\subsection{Arquitectura de red e infraestructura}

La topología de despliegue se diseña bajo la restricción de presupuesto acotado (RNF-INF-01), sobre servicios gestionados de nivel gratuito. Las aplicaciones cliente —la app móvil, distribuida como paquete instalable, y el panel web, publicado como sitio estático— se comunican con la API REST exclusivamente sobre HTTPS, adjuntando el token de sesión en cada solicitud. La API se ejecuta dentro de un contenedor sobre una plataforma como servicio (PaaS) y se conecta a una base de datos PostgreSQL gestionada y a un almacenamiento de objetos para las imágenes de los paquetes. El servicio de ruteo es opcional: ante su ausencia, el sistema degrada de forma automática a una estimación de distancias propia (RNF-AVL-01).

El flujo de integración y despliegue continuo se apoya en el control de versiones: cada publicación de cambios dispara el despliegue automático del backend y del panel web, mientras que el paquete de la aplicación móvil se genera mediante el servicio de compilación correspondiente. Esta elección mantiene la portabilidad del sistema (RNF-INF-02): migrar de proveedor implica cambiar dónde se ejecuta el mismo contenedor, no el código de la aplicación.

La Figura~\ref{fig:despliegue} sintetiza la topología de despliegue y el flujo de integración y despliegue continuo.

\begin{figure}[ht]
    \centering\color{nuevo}
    \begin{tikzpicture}
        \node[c4ext, text width=2.4cm] (dev) at (-6,1.6)  {Desarrollo local (contenedores)};
        \node[c4ext] (gh) at (-6,-1) {Repositorio\\ (GitHub)};
        \node[c4container] (ren) at (1.2,1.7) {API\\ (contenedor)};
        \node[c4container] (ver) at (4.4,1.7) {Panel web\\ (sitio estático)};
        \node[c4db] (sup) at (2.7,-0.6) {Base de datos\\ + almacenamiento};
        \node[c4ext] (eas) at (-2.6,-3.5) {Servicio de\\ compilación};
        \node[c4ext] (and) at (1.2,-3.5)  {Teléfono\\ (paquete móvil)};
        \draw[c4arrow] (dev) -- (gh) node[c4lbl] {publica};
        \draw[c4arrow] (gh) -- (ren) node[c4lbl, pos=0.68] {despliegue};
        \draw[c4arrow] (gh) -- (ver) node[c4lbl, pos=0.48] {despliegue};
        \draw[c4arrow] (gh) -- (eas) node[c4lbl] {compila};
        \draw[c4arrow] (eas) -- (and) node[c4lbl] {paquete};
        \draw[c4arrow] (ren) -- (sup);
        \draw[c4arrow, dashed] (ver) -- (ren) node[c4lbl] {API};
        \draw[c4arrow] (and) -- (ren) node[c4lbl, pos=0.55] {HTTPS};
        \node[boundary, fit=(ren)(ver)(sup), inner sep=5mm, label=above:{\footnotesize Nube --- nivel gratuito}] {};
    \end{tikzpicture}
    \caption{Diagrama de despliegue: topología de infraestructura e integración continua}
    \label{fig:despliegue}
\end{figure}

% ═════════════════════════════════════════════════════════════════════════════
\section{Tecnologías utilizadas}

La selección de tecnologías responde a los requerimientos no funcionales y a las decisiones de diseño fundamentadas en los antecedentes. La Tabla~\ref{tab:stack} resume el \emph{stack} adoptado y su justificación.

\begin{longtable}{p{3.1cm}p{4.4cm}p{6cm}}
    \caption{Tecnologías seleccionadas y su justificación}
    \label{tab:stack}\\
    \toprule
    \textbf{Capa} & \textbf{Tecnología} & \textbf{Justificación} \\
    \midrule
    \endfirsthead
    \multicolumn{3}{c}{\tablename\ \thetable\ -- \textit{continuación}} \\
    \toprule
    \textbf{Capa} & \textbf{Tecnología} & \textbf{Justificación} \\
    \midrule
    \endhead
    \bottomrule
    \endlastfoot

    Backend & Python con FastAPI (asíncrono) & Concurrencia nativa para operaciones de entrada/salida, documentación de API autogenerada (RNF-MNT-03) y validación de datos integrada (RNF-SEC-04). \\
    Persistencia & PostgreSQL con extensión geoespacial & Soporte nativo de tipos y consultas geográficas, necesario para la compatibilidad geográfica del \emph{matching}; disponible en nivel gratuito (RNF-INF-01). \\
    Acceso a datos & Mapeo objeto-relacional asíncrono & Aísla la lógica de negocio del motor de base de datos y facilita las pruebas (patrón repositorio, RNF-MNT-01). \\
    Autenticación & Tokens de acceso y refresco con \emph{hashing} de contraseñas & Sesiones seguras sin estado y almacenamiento resistente de credenciales (RNF-SEC-01). \\
    Ruteo y distancias & Servicio de ruteo por calles con alternativa propia & Cálculo de desvío real y de distancias para el \emph{matching} y el CO\textsubscript{2}, con degradación automática (RNF-AVL-01). \\
    Modelo de IA & Red convolucional preentrenada con aprendizaje por transferencia & Clasificación volumétrica de la carga con un conjunto de datos acotado, entrenada por el equipo (RNF-UADE-01). \\
    Aplicación móvil & React Native con Expo & Desarrollo multiplataforma desde una única base de código, con distribución de paquetes instalables. \\
    Panel web & Aplicación de página única & Interfaz de escritorio para PyMEs y administración (tablas densas y gráficos) que consume la misma API. \\
    Almacenamiento de imágenes & Almacenamiento de objetos con URLs firmadas & Acceso restringido y temporal a las fotos de los paquetes (RNF-SEC-03). \\
    Infraestructura & Contenedores sobre plataforma como servicio & Portabilidad (RNF-INF-02) y despliegue sin administración de servidores, en nivel gratuito (RNF-INF-01). \\
    Observabilidad & Registro estructurado de eventos & Trazabilidad de la operación (RNF-OBS-01). \\
    Pruebas e integración & Marco de pruebas automatizadas e integración continua & Cobertura de los módulos críticos y verificación en cada cambio (RNF-MNT-02). \\
\end{longtable}

En el núcleo del componente de inteligencia artificial se emplea aprendizaje por transferencia sobre una red convolucional preentrenada: en lugar de entrenar un modelo desde cero —inviable con el volumen de datos disponible—, se reutiliza una arquitectura con capacidades visuales genéricas y se la adapta al dominio del problema, cumpliendo el requisito de un modelo entrenado por el equipo (RNF-UADE-01) y la restricción de inferencia local sin servicios de terceros.

% ═════════════════════════════════════════════════════════════════════════════
\section{Modelo de datos}

El modelo de datos se organiza en torno a las entidades del dominio. La Tabla~\ref{tab:entidades} describe las principales; las relaciones más relevantes vinculan a un usuario con su perfil de transportista, a este con sus trayectos publicados, y a cada envío con su cliente, su transportista asignado, su clasificación de carga y sus eventos de estado.

\begin{longtable}{p{3.2cm}p{10.4cm}}
    \caption{Entidades principales del modelo de datos}
    \label{tab:entidades}\\
    \toprule
    \textbf{Entidad} & \textbf{Descripción} \\
    \midrule
    \endfirsthead
    \multicolumn{2}{c}{\tablename\ \thetable\ -- \textit{continuación}} \\
    \toprule
    \textbf{Entidad} & \textbf{Descripción} \\
    \midrule
    \endhead
    \bottomrule
    \endlastfoot

    Usuario & Datos de contacto y credenciales; tipo (particular o comercial) y roles (cliente, transportista). \\
    Transportista & Perfil asociado a un usuario: tipo de vehículo, licencia, patente, estado de validación, reputación y capacidad. \\
    Trayecto & Recorrido habitual (colaborativo) o ventana de disponibilidad (dedicada), con origen, destino, geometría y ventana horaria. \\
    Envío & Solicitud de transporte: origen, destino, modalidad, categoría de carga, estado, estado de pago, precio, ventana horaria y CO\textsubscript{2} ahorrado. \\
    Evento de envío & Registro de auditoría de cada transición de estado del envío, con marca temporal y ubicación. \\
    Clasificación & Resultado de la clasificación de la carga: imagen, categoría predicha, confianza, aceptación y categoría manual eventual. \\
    Traza GPS & Posiciones publicadas por el transportista durante un envío, para el seguimiento en tiempo real. \\
    Calificación & Puntuación y comentario asociados a un envío entregado. \\
    Pesos de asignación & Parámetros configurables del algoritmo de \emph{matching}, editables por administración. \\
    Organización & PyME (comercio y/o flota) con sus miembros, su flota de transportistas vinculados y sus envíos. \\
    \bottomrule
\end{longtable}

Los atributos geográficos (orígenes, destinos, geometrías de trayecto y trazas) se modelan con tipos espaciales e índices específicos, que sustentan las consultas de proximidad del \emph{matching}. El ciclo de vida del envío se gobierna mediante una máquina de estados (pendiente $\rightarrow$ asignado $\rightarrow$ retirado $\rightarrow$ en tránsito $\rightarrow$ entregado, con transición a cancelado), donde cada cambio válido se registra como un evento de auditoría (RF-SHP-05).

La Figura~\ref{fig:der} presenta el diagrama entidad-relación con las entidades principales y sus cardinalidades.

\begin{figure}[ht]
    \centering\color{nuevo}
    \resizebox{0.96\textwidth}{!}{%
    \begin{tikzpicture}[
        rel/.style={-, thick, black!70},
        card/.style={font=\scriptsize, fill=white, inner sep=1pt},
    ]
        \node[entity, text width=2.5cm] (usr)  at (-5.6,2.6)  {\textbf{Usuario}\nodepart{second}email, tipo, roles};
        \node[entity, text width=2.6cm] (tra)  at (-5.6,-1)   {\textbf{Transportista}\nodepart{second}vehículo, patente, estado, reputación};
        \node[entity, text width=2.6cm] (tray) at (-5.6,-4.7) {\textbf{Trayecto}\nodepart{second}tipo, origen, destino, ventana};
        \node[entity, text width=3cm]   (env)  at (0,-1)      {\textbf{Envío}\nodepart{second}origen, destino, modalidad, categoría, estado, precio, CO\textsubscript{2}};
        \node[entity, text width=2.5cm] (cla)  at (0,2.7)     {\textbf{Clasificación}\nodepart{second}categoría, confianza};
        \node[entity, text width=2.4cm] (evt)  at (5.2,2.7)   {\textbf{Evento de envío}\nodepart{second}estado, marca temporal};
        \node[entity, text width=2.4cm] (cal)  at (5.2,-1)    {\textbf{Calificación}\nodepart{second}estrellas, comentario};
        \node[entity, text width=2.4cm] (gps)  at (5.2,-4.7)  {\textbf{Traza GPS}\nodepart{second}posición, marca temporal};
        \node[entity, text width=2.5cm] (org)  at (0,-5)      {\textbf{Organización}\nodepart{second}nombre, CUIT, tipo};
        \draw[rel] (usr)  -- (tra)  node[card,pos=0.12]{1} node[card,pos=0.9]{0..1};
        \draw[rel] (tra)  -- (tray) node[card,pos=0.12]{1} node[card,pos=0.9]{N};
        \draw[rel] (usr)  -- (env)  node[card,pos=0.12]{1} node[card,pos=0.9]{N};
        \draw[rel] (tra)  -- (env)  node[card,pos=0.12]{1} node[card,pos=0.9]{N};
        \draw[rel] (env)  -- (cla)  node[card,pos=0.12]{1} node[card,pos=0.9]{0..1};
        \draw[rel] (env)  -- (evt)  node[card,pos=0.12]{1} node[card,pos=0.9]{N};
        \draw[rel] (env)  -- (cal)  node[card,pos=0.12]{1} node[card,pos=0.9]{N};
        \draw[rel] (env)  -- (gps)  node[card,pos=0.12]{1} node[card,pos=0.9]{N};
        \draw[rel] (org)  -- (env)  node[card,pos=0.12]{1} node[card,pos=0.9]{N};
        \draw[rel] (org)  -- (tra)  node[card,pos=0.12]{1} node[card,pos=0.9]{N};
    \end{tikzpicture}%
    }
    \caption{Diagrama entidad-relación del modelo de datos}
    \label{fig:der}
\end{figure}

% ═════════════════════════════════════════════════════════════════════════════
\section{Validación del sistema}

La validación combina pruebas automatizadas del software y evaluación cuantitativa del modelo de inteligencia artificial.

En el backend, los módulos críticos —autenticación, envíos y asignación— se cubren con pruebas automatizadas, con el objetivo de una cobertura de al menos el 60\% (RNF-MNT-02). El diseño favorece esta verificación: los cálculos centrales del \emph{matching} y del CO\textsubscript{2} se implementan como funciones puras y determinísticas, lo que permite comprobar de forma directa la tabla de compatibilidad vehículo-carga, los filtros duros (descarte de la modalidad colaborativa para cargas voluminosas, de la movilidad a pie o en bicicleta en trayectos largos y de los desvíos por encima del máximo admitido), el ordenamiento por proximidad y el cálculo del ahorro de emisiones con coordenadas reales del AMBA.

El modelo de clasificación de carga se evalúa exclusivamente sobre el conjunto de prueba, reservado hasta la instancia final. Se reportan la exactitud global, la precisión, la exhaustividad y la medida F por categoría, y la matriz de confusión, prestando atención a las confusiones entre categorías adyacentes. Complementariamente, se realiza un análisis de sesgos que agrupa el conjunto de prueba según las condiciones de captura —iluminación, ángulo y fondo— y según la presencia de un objeto de referencia, calculando la exactitud por grupo para detectar y documentar las condiciones en las que el desempeño se degrada (RNF-UADE-01). Este análisis alimenta tanto las mitigaciones del modelo como las recomendaciones de uso en la interfaz.
}% <<< FIN CONTENIDO NUEVO (entrega 50%) >>>
